\scanpage{52}
homéomorphismes universels\ldots

Les réflexions précédentes suggèrent aussi des énoncés comme les suivants : Pour un schéma de base $S$ loc. noeth. donné, le foncteur $X\longrightarrow X_{\mathrm{et}}$, allant de la catégorie \deleted{réduits} des schémas\inserted{ réduits} loc. de prés. finie sur $S$, vers la 2-catégorie des \unsure{topos} au dessus de $X_{\mathrm{et}}$, est \deleted{pleinement} 1-fidèle (deux hom. $f,g:X\unsure{=|\rightrightarrows}Y$ tels que les homomorphismes\inserted{ $f_{\mathrm{et}},g_{\mathrm{et}}$} de topos $X_{\mathrm{et}}\rightrightarrows Y_{\mathrm{et}}$ au-dessus de $S_{\mathrm{et}}$ soient \unsure{ég|$fg$} isomorphes, sont égaux) et $\unsure{=|\overline m}$ peut-être \uline{pleinement} \uline{fidèle}, quand on passe à la catégorie des fractions de $(\mathrm{Sch}_{\unsure{l.t.f.}})/S$ obtenue en rendant inv. les homéom. universels\ldots Expérimenter \deleted{\unsure{p}} ceci p.ex. pour les autom d'une courbe alg. propre sur une ext. finie de $\mathbb{Q}$, en \inserted{re}trouvant le ``fait'' que tout autom ext. de $E_K$ ($K$ le corps des fonctions de $X$) qui respecte la structure à \unsure{lacets|Qcts} et qui commute à l'action de $\Gamma$, provient d'un autom de $X$
\margin{$S$ de car. $0$ ?}
\margin{51}
